\section*{Lampiran B: Source Code Program Python AI (Random Forest Regressor)}
\phantomsection
\addcontentsline{toc}{section}{Lampiran B: Source Code Program Python AI}

Kode berikut merupakan program Python yang dijalankan pada Visual Studio Code.
Program membaca data historis suhu dari ThingSpeak \textit{public channel},
melatih model \textit{Random Forest Regressor}, dan mengirim hasil prediksi
suhu 1 menit ke depan beserta status klasifikasi kembali ke ThingSpeak.
Pengiriman AI dibatasi minimal 45 detik antar pengiriman untuk menghindari
\textit{rate limit} ThingSpeak.

\begin{lstlisting}[style=pythonstyle,
    caption={Program Python AI Random Forest Regressor untuk Prediksi Suhu 1 Menit},
    label={lst:python}]
import time
import requests
import numpy as np
import pandas as pd

from sklearn.ensemble import RandomForestRegressor
from sklearn.metrics import mean_absolute_error, mean_squared_error, r2_score

# =========================================================
# KONFIGURASI THINGSPEAK
# =========================================================
CHANNEL_ID    = "3413573"
READ_API_KEY  = ""               # channel public
WRITE_API_KEY = "ISI_WRITE_API_KEY_ANDA"

READ_URL  = f"https://api.thingspeak.com/channels/{CHANNEL_ID}/feeds.json"
WRITE_URL = "https://api.thingspeak.com/update"

# =========================================================
# PARAMETER SISTEM
# =========================================================
DATA_INTERVAL_SECOND      = 15   # ESP32 kirim tiap 15 detik
PREDICTION_HORIZON_SECOND = 60   # prediksi 1 menit ke depan
PREDICTION_STEPS = int(PREDICTION_HORIZON_SECOND / DATA_INTERVAL_SECOND)
# PREDICTION_STEPS = 4

AI_READ_INTERVAL_SECOND  = 10   # baca data tiap 10 detik
AI_WRITE_INTERVAL_SECOND = 45   # tulis prediksi minimal tiap 45 detik
MIN_FEATURE_DATA = 5            # minimal data fitur untuk training

# =========================================================
# KLASIFIKASI SUHU: 0=Dingin, 1=Normal, 2=Panas
# =========================================================
def classify_temperature(temp: float) -> int:
    if temp < 25.0:   return 0
    elif temp <= 30.0: return 1
    else:              return 2

def label_text(label: int) -> str:
    return {0: "DINGIN", 1: "NORMAL", 2: "PANAS"}.get(label, "TIDAK DIKETAHUI")

# =========================================================
# MEMBACA DATA DARI THINGSPEAK PUBLIC CHANNEL
# =========================================================
def read_temperature_data(results: int = 200):
    try:
        response = requests.get(READ_URL, params={"results": results}, timeout=10)
    except requests.exceptions.RequestException as e:
        print("Gagal koneksi ke ThingSpeak:", e)
        return None

    if response.status_code != 200 or response.text.strip() == "-1":
        print("Gagal membaca data. Status:", response.status_code)
        return None

    try:
        feeds = response.json().get("feeds", [])
    except ValueError:
        print("Response ThingSpeak bukan JSON.")
        return None

    rows = []
    for feed in feeds:
        val = feed.get("field1")
        if val is not None:
            try:
                rows.append({"entry_id":    feed.get("entry_id"),
                             "created_at":  feed.get("created_at"),
                             "temperature": float(val)})
            except ValueError:
                pass

    if not rows:
        print("Belum ada data suhu valid pada Field 1 ThingSpeak")
        return None

    return pd.DataFrame(rows).dropna().reset_index(drop=True)

# =========================================================
# MEMBUAT FITUR TIME SERIES
# =========================================================
def create_time_series_features(df: pd.DataFrame):
    df = df.copy()
    df["temp_lag_1"]   = df["temperature"].shift(1)
    df["temp_lag_2"]   = df["temperature"].shift(2)
    df["temp_lag_3"]   = df["temperature"].shift(3)
    df["temp_change"]  = df["temperature"] - df["temp_lag_1"]
    df["moving_avg_3"] = df["temperature"].rolling(window=3).mean()
    df["moving_avg_5"] = df["temperature"].rolling(window=5).mean()
    df["target_future"]= df["temperature"].shift(-PREDICTION_STEPS)
    return df.dropna().reset_index(drop=True)

# =========================================================
# TRAINING RANDOM FOREST REGRESSOR
# =========================================================
FEATURE_COLS = ["temperature","temp_lag_1","temp_lag_2","temp_lag_3",
                "temp_change","moving_avg_3","moving_avg_5"]

def train_random_forest_regressor(df_feature: pd.DataFrame):
    if len(df_feature) < MIN_FEATURE_DATA:
        print(f"Data belum cukup. Minimal: {MIN_FEATURE_DATA}")
        return None

    X, y     = df_feature[FEATURE_COLS], df_feature["target_future"]
    split    = int(len(df_feature) * 0.8)

    if split < 2:
        X_train, y_train = X, y
        X_test,  y_test  = X, y
    else:
        X_train, X_test = X.iloc[:split], X.iloc[split:]
        y_train, y_test = y.iloc[:split], y.iloc[split:]

    model = RandomForestRegressor(n_estimators=100, random_state=42, max_depth=6)
    model.fit(X_train, y_train)

    if len(X_test) > 0:
        y_pred = model.predict(X_test)
        mae  = mean_absolute_error(y_test, y_pred)
        rmse = np.sqrt(mean_squared_error(y_test, y_pred))
        r2   = r2_score(y_test, y_pred) if len(y_test) > 1 else 0.0
        print(f"MAE={mae:.3f} C  RMSE={rmse:.3f} C  R2={r2:.3f}")
    return model

# =========================================================
# INPUT TERBARU UNTUK PREDIKSI
# =========================================================
def create_latest_input(df: pd.DataFrame):
    if len(df) < 5:
        print("Minimal 5 data suhu untuk input prediksi.")
        return None
    t = df["temperature"]
    return pd.DataFrame([{
        "temperature":  t.iloc[-1],
        "temp_lag_1":   t.iloc[-2],
        "temp_lag_2":   t.iloc[-3],
        "temp_lag_3":   t.iloc[-4],
        "temp_change":  t.iloc[-1] - t.iloc[-2],
        "moving_avg_3": t.iloc[-3:].mean(),
        "moving_avg_5": t.iloc[-5:].mean()
    }])

# =========================================================
# MENGIRIM HASIL AI KE THINGSPEAK
# Field 2 = prediksi suhu, Field 3 = status
# =========================================================
def send_prediction_to_thingspeak(pred_temp: float, status: int):
    params = {"api_key": WRITE_API_KEY,
              "field2":  round(pred_temp, 2),
              "field3":  status}
    try:
        resp = requests.get(WRITE_URL, params=params, timeout=10)
        if resp.status_code == 200 and resp.text.strip() != "0":
            print("Prediksi AI terkirim. Entry ID:", resp.text)
            return True
        print("Gagal kirim. Response:", resp.text)
        return False
    except requests.exceptions.RequestException as e:
        print("Error kirim:", e)
        return False

# =========================================================
# PROGRAM UTAMA
# =========================================================
def main():
    print("Program AI Random Forest - Prediksi Suhu 1 Menit ke Depan")
    print(f"Interval sensor   : {DATA_INTERVAL_SECOND} detik")
    print(f"Horizon prediksi  : {PREDICTION_HORIZON_SECOND} detik ({PREDICTION_STEPS} langkah)")
    print(f"Interval baca AI  : {AI_READ_INTERVAL_SECOND} detik")
    print(f"Interval tulis AI : {AI_WRITE_INTERVAL_SECOND} detik")
    print()

    last_write = 0

    while True:
        df = read_temperature_data(results=200)
        if df is not None:
            print(f"Data terbaca: {len(df)} | Entry: {df['entry_id'].iloc[-1]}"
                  f" | Suhu: {df['temperature'].iloc[-1]:.2f} C")

            df_feat = create_time_series_features(df)
            if len(df_feat) < MIN_FEATURE_DATA:
                print(f"Data fitur kurang ({len(df_feat)}/{MIN_FEATURE_DATA})")
                time.sleep(AI_READ_INTERVAL_SECOND)
                continue

            model = train_random_forest_regressor(df_feat)
            if model is not None:
                inp = create_latest_input(df)
                if inp is not None:
                    pred_temp   = model.predict(inp)[0]
                    pred_status = classify_temperature(pred_temp)
                    print(f"Aktual: {df['temperature'].iloc[-1]:.2f} C"
                          f" | Prediksi 1 mnt: {pred_temp:.2f} C"
                          f" | Status: {label_text(pred_status)}")

                    now = time.time()
                    if now - last_write >= AI_WRITE_INTERVAL_SECOND:
                        if send_prediction_to_thingspeak(pred_temp, pred_status):
                            last_write = now
                    else:
                        sisa = AI_WRITE_INTERVAL_SECOND - (now - last_write)
                        print(f"Tunggu {sisa:.0f} detik sebelum kirim AI.")

        time.sleep(AI_READ_INTERVAL_SECOND)

if __name__ == "__main__":
    main()
\end{lstlisting}
